
\documentclass{article}
\usepackage{colortbl}
\usepackage{makecell}
\usepackage{multirow}
\usepackage{supertabular}

\begin{document}

\newcounter{utterance}

\centering \large Interaction Transcript for game `textmapworld', experiment `large', episode 25 with reka{-}flash{-}3.1{-}t0.0.
\vspace{24pt}

{ \footnotesize  \setcounter{utterance}{1}
\setlength{\tabcolsep}{0pt}
\begin{supertabular}{c@{$\;$}|p{.15\linewidth}@{}p{.15\linewidth}p{.15\linewidth}p{.15\linewidth}p{.15\linewidth}p{.15\linewidth}}
    \# & $\;$A & \multicolumn{4}{c}{Game Master} & $\;\:$B\\
    \hline

    \theutterance \stepcounter{utterance}  
    & & \multicolumn{4}{p{0.6\linewidth}}{
        \cellcolor[rgb]{0.9,0.9,0.9}{
            \makecell[{{p{\linewidth}}}]{
                \texttt{\tiny{[P1$\langle$GM]}}
                \texttt{Please help me with the following task. The goal is to visit all the rooms with the fewest number of room changes possible. In each room, you need to decide the direction to go in. Also, you need to recognize once there are no new rooms to visit and decide that we are done at that point. Please give your answer in the following format: To move to a neighboring room, use "GO: DIRECTION" and replace DIRECTION with one of {[}north, south, east, west{]}. To stop the exploration, answer with "DONE" instead. Omit any other text.} \\
\texttt{Here is an example:} \\
\texttt{You are in the Kitchen. Currently available directions: south, west. What is your next instruction?} \\
\texttt{GO: west} \\
\texttt{You have made a step and entered a Lobby. Currently available directions: east, north. What is your next instruction?} \\
\texttt{GO: north} \\
\texttt{...} \\
\texttt{You have made a step and entered a Bedroom. Currently available directions: south. What is your next instruction?} \\
\texttt{DONE} \\
\texttt{Let us start. You are in the Pantry. Currently available directions: north, south. What is your next instruction?} \\
            }
        }
    }
    & \\ \\

    \theutterance \stepcounter{utterance}  
    & \multicolumn{4}{p{0.6\linewidth}}{
        \cellcolor[rgb]{0.9,0.9,0.9}{
            \makecell[{{p{\linewidth}}}]{
                \texttt{\tiny{[P1$\rangle$GM]}}
                \texttt{<reasoning>} \\
\texttt{Let's tackle this problem. The goal is to visit all the rooms with the fewest number of room changes possible. So, I need to figure out the optimal path that covers all rooms in the least steps. But how do I know the layout of the rooms? The user hasn't provided a map, so I must be working with the information given step by step as I move through the rooms.} \\
\\ 
\texttt{First, the starting point is the Pantry with available directions north and south. The challenge is to explore each room once, but since the problem mentions "visit all the rooms," I assume there are multiple rooms connected in some way. The key here is to minimize the number of room changes, which I think refers to the number of moves (i.e., steps) taken. So, it's similar to finding the shortest path that covers all rooms, which sounds like the Traveling Salesman Problem (TSP). But TSP is NP{-}hard, so maybe there's a simpler approach here given the dynamic nature of the exploration.} \\
\\ 
\texttt{Wait, but since the rooms are connected in a graph structure, perhaps it's a maze{-}like layout. The problem requires making decisions at each step based on available directions. The example given shows moving west, then north, etc., until all rooms are visited. So, the strategy might involve backtracking when necessary to ensure all rooms are covered, but trying to do so with minimal steps.} \\
\\ 
\texttt{But how do I track which rooms have been visited? Since the user hasn't provided a map, I need to keep track of visited rooms as I go. Each time I enter a room, I note it down. Then, when choosing directions, I should prioritize unvisited rooms first. If all adjacent rooms have been visited, then I have to backtrack. This is similar to depth{-}first search (DFS) but with the goal of minimizing steps. However, DFS can lead to longer paths, so maybe a breadth{-}first search (BFS) approach would be better? But BFS is for finding the shortest path between two points, not necessarily covering all nodes.} \\
\\ 
\texttt{Alternatively, since the problem is about visiting all rooms with the fewest room changes (steps), it's equivalent to finding the shortest possible route that visits every room at least once. This is exactly the Hamiltonian path problem, which is also NP{-}complete. But given that the actual layout isn't known, perhaps the solution is to explore greedily, always choosing an unvisited room if possible} \\
            }
        }
    }
    & & \\ \\

    \theutterance \stepcounter{utterance}  
    & & & \multicolumn{2}{p{0.3\linewidth}}{
        \cellcolor[rgb]{0.9,0.9,0.9}{
            \makecell[{{p{\linewidth}}}]{
                \texttt{\tiny{[GM$|$GM]}}
                \texttt{abort game} \\
            }
        }
    }
    & & \\ \\

\end{supertabular}
}

\end{document}
